\documentclass[11pt]{article}
\usepackage[utf8]{inputenc}
\usepackage[T1]{fontenc}
\usepackage{amsmath,amssymb,amsthm}
\usepackage{geometry}
\usepackage{longtable,array,booktabs}
\usepackage{enumitem}
\usepackage{hyperref}
\geometry{margin=1in}
\hypersetup{colorlinks=true,linkcolor=blue,citecolor=blue,urlcolor=blue}

\newtheorem{theorem}{Theorem}[section]
\newtheorem{lemma}[theorem]{Lemma}
\newtheorem{proposition}[theorem]{Proposition}
\newtheorem{corollary}[theorem]{Corollary}
\theoremstyle{definition}
\newtheorem{definition}[theorem]{Definition}
\newtheorem{assumption}[theorem]{Assumption}
\newtheorem{remark}[theorem]{Remark}

\newcommand{\C}{\mathbb{C}}
\newcommand{\R}{\mathbb{R}}
\newcommand{\N}{\mathbb{N}}
\newcommand{\Z}{\mathbb{Z}}
\newcommand{\dd}{\,\mathrm{d}}
\newcommand{\Capacity}{\operatorname{Cap}}
\newcommand{\dist}{\operatorname{dist}}
\newcommand{\spec}{\operatorname{spec}}
\newcommand{\erfc}{\operatorname{erfc}}
\newcommand{\ord}{\operatorname{ord}}
\newcommand{\Res}{\operatorname{Res}}
\newcommand{\diag}{\operatorname{diag}}
\newcommand{\ML}[2]{E_{#1,#2}}
\newcommand{\Lc}{\mathcal{L}}
\newcommand{\Dc}{\mathcal{D}}
\newcommand{\Ig}{I^{\mu}}
\newcommand{\Pade}[1]{[#1/#1]}
\newcommand{\Source}[1]{\par\smallskip\noindent\textup{[Sources: #1.]}}

\title{A Unified Exposition of the Fractional Riccati--Müntz--Padé Theory}
\author{arb4j mathematical documentation synthesis}
\date{\today}

\begin{document}
\maketitle

\begin{abstract}
This article reconciles the mathematical documentation in \texttt{docs/} into one canonical account.
The object is the constant-coefficient fractional Riccati equation,
its Puiseux--Müntz solution,
the Müntz--Tau recurrence,
Cole--Hopf meromorphy,
diagonal Padé continuation in the variable $z=t^{\mu}$,
Hankel and orthogonal-polynomial constructions,
Chebyshev--Wheeler computation,
Jacobi and Toda interpretations,
Riemann--Hilbert and Stahl convergence theory,
and the rough-Heston pricing application.
The exposition fixes one notation and states each major result once.
Where drafts use different index conventions or hypotheses,
the canonical version below records the translation and the sharper condition.
\end{abstract}

\tableofcontents

\section*{Notation}
\addcontentsline{toc}{section}{Notation}

Throughout,
$0<\mu\le 1$ is the fractional order.
In rough Heston, $\mu=H+\tfrac12$ with $H\in(0,\tfrac12]$.
The physical time is $t\ge0$.
The Müntz variable is
\[
        z=t^{\mu}.
\]
The scalar fractional Riccati equation is written
\begin{equation}\label{eq:canonical-riccati}
        D^{\mu}y(t)=c_0+c_1y(t)+c_2y(t)^2,
        \qquad y(0)=0,
\end{equation}
where $D^{\mu}$ is the Caputo derivative.
Equivalently,
\begin{equation}\label{eq:canonical-volterra}
        y(t)=I^{\mu}\{c_0+c_1y+c_2y^2\}(t).
\end{equation}
For rough Heston the constants become polynomials in the Fourier variable.
The correspondence is
\[
        c_0=P(u),\qquad c_1=Q(u),\qquad c_2=R(u).
\]
The pricing Fourier variable is denoted $v$.
The symbol $u$ is reserved for the Riccati parameter in the structural theory;
when the pricing transform is used we set $u=v$ or $u=-iw$ according to the contour convention.

The Puiseux coefficients are $a_k$ and the canonical series is
\begin{equation}\label{eq:ak-series}
        y(t)=\sum_{k\ge1}a_k t^{k\mu}
        =\sum_{k\ge1}a_k z^k.
\end{equation}
Some source files index the same sequence as $a(k)$ with $a(k)$ multiplying $t^{(k+1)\mu}$.
The translation is $a(k)=a_{k+1}$.
For the cumulant generating function,
\begin{equation}\label{eq:dk-series}
        \Phi(v,T)=\sum_{k\ge0} d_k(v) T^{(k+1)\mu}
        =T^{\mu}\sum_{k\ge0} d_k(v) z^k
        \quad (z=T^{\mu}).
\end{equation}
The $d_k$ are not renamed to $a_k$ because they include the time-integration weights and Heston coefficients.

The diagonal Padé approximant in $z$ is
\begin{equation}\label{eq:pade-symbol}
        \Phi_M(z;v)=\frac{\Phi^{\mathrm{num}}_M(z;v)}{\Phi^{\mathrm{den}}_M(z;v)}.
\end{equation}
For a generic scalar function $g(z)=\sum_{k\ge0}m_kz^k$ the Hankel matrix is
\[
        H_M=(m_{i+j})_{0\le i,j<M},
        \qquad \tau_M=\det H_M.
\]
The monic orthogonal polynomials for the moment functional are $p_n$,
the recurrence coefficients are $\alpha_n,\beta_n$,
and the squared norms are
\[
        h_n=\Lc[p_n^2]=\Lc[x^np_n]=\frac{\tau_{n+1}}{\tau_n}.
\]
The diagonal denominator is the reciprocal polynomial
\[
        Q_M(z)=z^M p_M(1/z).
\]
The Padé poles are denoted $\eta_j$.
Residues at those poles are denoted $r_j$;
this avoids conflict with the Riccati coefficients $c_0,c_1,c_2$.
The spectral measure attached to a Cauchy transform of $g$ is denoted $\nu_g$,
never by the fractional-order symbol $\mu$.
The pricing contour variable after the Lewis rotation is $w=c+i\xi$.

The source documents use four principal normalisations.
First,
\texttt{frmp.tex} uses $a_k$ for Riccati coefficients and $d_k$ for cgf coefficients;
that convention is authoritative here.
Second,
\texttt{ThePadeMuntzSpectralTauMethodForConstantCoeffecientFractionalRiccations.tex}
uses $g(z)=y(z^{1/\mu})$ and writes the Müntz--Caputo operator as $\Dc$.
Third,
\texttt{FractionalRiccatiSolutionMethodology.tex} writes $S(n,z)$ for generating polynomials;
we treat them as reciprocal orthogonal-polynomial data.
Fourth,
\texttt{TheRiemannHilbertViewOfPadeApproximationAndGlobalAnalyticContinuation.tex}
uses $R_M=P_M/Q_M$ for the rational approximant;
we use the same notation in the general Padé sections and reserve $\Phi_M$ for the cgf.

\section{Fractional calculus preliminaries}

\begin{definition}[Riemann--Liouville integral and Caputo derivative]\label{def:frac-ops}
For $\mu>0$ the Riemann--Liouville integral is
\begin{equation}\label{eq:rl-integral}
        (I^{\mu}f)(t)=\frac1{\Gamma(\mu)}\int_0^t (t-s)^{\mu-1}f(s)\dd s.
\end{equation}
For $0<\mu\le1$ and absolutely continuous $f$,
\begin{equation}\label{eq:caputo}
        (D^{\mu}f)(t)=I^{1-\mu}f'(t).
\end{equation}
If $f(0)=0$, then $D^{\mu}I^{\mu}f=f$ and $I^{\mu}D^{\mu}f=f$ on the natural Volterra class.
\Source{frmp.tex Def. \texttt{def:RL}; ThePadeMuntz\ldots Defs. \texttt{def:RL}, \texttt{def:Caputo}; RH-view Def. ``Riemann--Liouville integral and Caputo derivative''.}
\end{definition}

\begin{lemma}[Power rule]\label{lem:power-rule}
For $\rho>-1$,
\begin{equation}\label{eq:power-int}
        I^{\mu}t^{\rho}
        =\frac{\Gamma(\rho+1)}{\Gamma(\rho+\mu+1)}t^{\rho+\mu}.
\end{equation}
For $\rho>\mu-1$,
\begin{equation}\label{eq:power-der}
        D^{\mu}t^{\rho}
        =\frac{\Gamma(\rho+1)}{\Gamma(\rho-\mu+1)}t^{\rho-\mu}.
\end{equation}
In particular,
\begin{equation}\label{eq:gamma-k}
        I^{\mu}t^{k\mu}
        =\gamma_{k+1}t^{(k+1)\mu},
        \qquad
        \gamma_{k+1}:=\frac{\Gamma(k\mu+1)}{\Gamma((k+1)\mu+1)}.
\end{equation}
\Source{frmp.tex Lem. \texttt{lem:power}; ThePadeMuntz\ldots Prop. \texttt{prop:powers}; RH-view Prop. ``Action on power functions''.}
\end{lemma}

\begin{proof}
Insert $s=tu$ in \eqref{eq:rl-integral}.
The remaining integral is the beta integral
$B(\rho+1,\mu)=\Gamma(\rho+1)\Gamma(\mu)/\Gamma(\rho+\mu+1)$.
The derivative formula follows by applying the preceding identity to $f'$ and simplifying gamma factors.
\end{proof}

\begin{remark}[Caputo versus Riemann--Liouville in the drafts]
Some drafts write the initial condition as $I^{1-\mu}y(0)=0$.
For the zero-initial-value Volterra solution used here this is equivalent to the Caputo formulation.
The Caputo notation is cleaner for the scalar IVP,
while the Riemann--Liouville integral is cleaner for deriving coefficients.
\end{remark}

\section{The Müntz lattice}

\begin{definition}[Müntz lattice]\label{def:muntz-lattice}
The Müntz lattice associated with $\mu$ is the linear span
\begin{equation}\label{eq:muntz-space}
        \mathcal{M}_{\mu}=\left\{\sum_{k\ge0}b_k z^k: z=t^{\mu}\right\}
\end{equation}
near $z=0$.
The subspace $z\mathcal{M}_{\mu}$ consists of series with zero constant term.
\Source{ThePadeMuntz\ldots Def. \texttt{def:muntz}; RH-view Def. ``Müntz lattice''.}
\end{definition}

\begin{proposition}[Algebra and fractional-calculus closure]\label{prop:muntz-closure}
The space $\mathcal{M}_{\mu}$ is closed under addition,
Cauchy products,
and multiplication by polynomials in $z$.
The integral $I^{\mu}$ maps $\mathcal{M}_{\mu}$ into $z\mathcal{M}_{\mu}$.
The Caputo derivative maps $z\mathcal{M}_{\mu}$ into $\mathcal{M}_{\mu}$.
\Source{ThePadeMuntz\ldots Prop. \texttt{prop:closure\_M}; RH-view Thm. ``Closure of the Müntz lattice''; FractionalRiccatiSolutionMethodology.tex Lem. \texttt{lem:conv}.}
\end{proposition}

\begin{proof}
Addition is coefficientwise.
Products follow from the Cauchy product
\[
        \left(\sum_{j\ge0}b_jz^j\right)
        \left(\sum_{\ell\ge0}c_\ell z^\ell\right)
        =\sum_{m\ge0}\left(\sum_{j+\ell=m}b_jc_\ell\right)z^m.
\]
The integral and derivative actions are exactly \eqref{eq:power-int} and \eqref{eq:power-der} on monomials.
\end{proof}

\begin{definition}[Müntz--Caputo operator]\label{def:muntz-caputo}
For $g(z)=y(t)$ with $z=t^{\mu}$ define
\begin{equation}\label{eq:muntz-caputo}
        (\Dc g)(z)=(D^{\mu}y)(t).
\end{equation}
Thus
\begin{equation}\label{eq:muntz-caputo-monomial}
        \Dc z^k=\frac{\Gamma(k\mu+1)}{\Gamma((k-1)\mu+1)}z^{k-1},
        \qquad k\ge1.
\end{equation}
\Source{ThePadeMuntz\ldots Def. \texttt{def:cD}; RH-view rough-Heston application.}
\end{definition}

\begin{lemma}[Leibniz identity on the Müntz algebra]\label{lem:muntz-leibniz}
On products needed for the Cole--Hopf transform,
\begin{equation}\label{eq:muntz-leibniz}
        \Dc(gw)=(\Dc g)w+g(\Dc w)
\end{equation}
whenever one factor is generated by the linear Müntz system used below.
\Source{ThePadeMuntz\ldots Lem. \texttt{lem:leibniz}; RH-view meromorphy synthesis.}
\end{lemma}

\begin{proof}
Both sides are compared coefficientwise on the common Müntz expansion.
The gamma weights match because the linear system fixes the same shift on every monomial.
This is the exact product rule used by the Cole--Hopf calculation;
it is not a general ordinary derivative rule in the $z$ variable.
\end{proof}

\section{The Müntz--Tau recurrence}

\begin{theorem}[Müntz--Tau recurrence]\label{thm:tau}
The Volterra equation \eqref{eq:canonical-volterra} has a unique formal solution
$y(t)=\sum_{k\ge1}a_kt^{k\mu}$.
The coefficients are
\begin{equation}\label{eq:a1-canonical}
        a_1=\frac{c_0}{\Gamma(\mu+1)},
\end{equation}
and for $k\ge2$,
\begin{equation}\label{eq:ak-canonical}
        a_k=\frac{\Gamma((k-1)\mu+1)}{\Gamma(k\mu+1)}
        \left(c_1a_{k-1}+c_2\sum_{j=1}^{k-2}a_ja_{k-1-j}\right).
\end{equation}
Equivalently,
\begin{equation}\label{eq:ak-shifted}
        a_{k+1}=\frac{\Gamma(k\mu+1)}{\Gamma((k+1)\mu+1)}
        \left(c_1a_k+c_2\sum_{j+\ell=k}a_ja_\ell\right),
        \qquad k\ge1,
\end{equation}
where the sum uses $j,\ell\ge1$ and is empty for $k=1$.
\Source{frmp.tex Thm. \texttt{thm:tau}; ThePadeMuntz\ldots Thm. \texttt{thm:riccati\_series}; RH-view Thm. ``Müntz--Tau Puiseux series solution''; TheFractionalRiccatiMuntzPadeTheorem.tm Lem. \texttt{lem:Muntz}. Statements differ only by the index shift $a(k)=a_{k+1}$.}
\end{theorem}

\begin{proof}
Substitute \eqref{eq:ak-series} into \eqref{eq:canonical-volterra}.
The square is
\[
        y(t)^2=\sum_{m\ge2}\left(\sum_{j=1}^{m-1}a_ja_{m-j}\right)t^{m\mu}.
\]
Apply Lemma~\ref{lem:power-rule} to the constant term,
the linear term,
and the quadratic Cauchy product.
The coefficient of $t^{\mu}$ gives \eqref{eq:a1-canonical}.
The coefficient of $t^{k\mu}$ for $k\ge2$ gives \eqref{eq:ak-canonical}.
\end{proof}

\begin{corollary}[Polynomial degree growth in the rough-Heston parameter]\label{cor:degree-growth}
If $c_0=P(u)$,
$c_1=Q(u)$,
and $c_2=R(u)$ are polynomials in $u$,
then each $a_k(u)$ is a polynomial in $u$.
For the rough-Heston symbols the degree grows at most linearly with $k$ after the initial quadratic Heston forcing is accounted for.
\Source{frmp.tex Lem. \texttt{lem:h}; FractionalRiccatiSolutionMethodology.tex Lem. \texttt{lem:degree}; TheFractionalRiccatiMuntzPadeTheorem.tm rough-Heston section.}
\end{corollary}

\begin{proof}
The recurrence uses only additions,
multiplications,
and scalar gamma factors.
Starting from the polynomial $a_1=P(u)/\Gamma(\mu+1)$,
induction proves polynomiality.
The displayed degree statement follows by adding degrees in the convolution and applying the Heston degree bounds for $P,Q,R$.
\end{proof}

\begin{definition}[Cumulant coefficients]\label{def:cgf-coefficients}
For rough Heston,
let $h(t,u)$ denote the fractional Riccati solution.
The cumulant exponent contains the time integral
\begin{equation}\label{eq:cgf-time-integral}
        \int_0^T h(s,u)\dd s
        =\sum_{k\ge1}\frac{a_k(u)}{k\mu+1}T^{k\mu+1}.
\end{equation}
After collecting the affine Heston terms,
the cgf coefficients $d_k(v)$ are the corresponding scalar multiples of $a_{k+1}(v)$ and the deterministic affine term.
\Source{frmp.tex Thm. \texttt{thm:cgf}; FractionalRiccatiSolutionMethodology.tex Thm. \texttt{thm:integ} and Cor. \texttt{cor:fracint}; TheFractionalRiccatiMuntzPadeTheorem.tm Cor. \texttt{cor:dk}.}
\end{definition}

\begin{theorem}[Cgf as a Müntz-lattice power series]\label{thm:cgf}
For admissible rough-Heston parameters and fixed $v$ near the pricing strip,
there are coefficients $d_k(v)$ such that
\begin{equation}\label{eq:cgf-main}
        \Phi(v,T)=\sum_{k\ge0}d_k(v)T^{(k+1)\mu}
\end{equation}
near $T=0$.
Each $d_k$ is a polynomial in $v$ determined by the Müntz--Tau recurrence and the affine time-integration weights.
\Source{frmp.tex Thm. \texttt{thm:cgf}; TheFractionalRiccatiMuntzPadeTheorem.tm rough-Heston application; FractionalRiccatiSolutionMethodology.tex Thm. \texttt{thm:roughHeston}.}
\end{theorem}

\begin{proof}[Proof sketch]
The El Euch--Rosenbaum fractional Riccati representation expresses the log transform through the solution $h$ and its time integral.
Insert Theorem~\ref{thm:tau}.
Termwise integration is valid inside the positive radius established in Theorem~\ref{thm:radius}.
The Heston affine constants add only polynomial contributions in $v$.
\end{proof}

\section{Radius of convergence and generic finiteness}

\begin{theorem}[Positive Puiseux radius]\label{thm:radius}
There exists $R>0$ such that the series \eqref{eq:ak-series} converges for $|z|<R$.
More concretely,
with
\begin{equation}\label{eq:T0}
        T_0^{\mu}=\min\left\{
        \frac{\Gamma(\mu+1)}{|c_0|+|c_1|+|c_2|+1},
        \frac{\Gamma(\mu+1)}{2(|c_1|+2|c_2|+1)}
        \right\},
\end{equation}
the Volterra operator is a contraction on the unit ball of $C[0,T_0]$,
and $R\ge T_0^{\mu}$.
For every $0<r<R$,
\begin{equation}\label{eq:cauchy-ak}
        |a_k|\le \frac{M(r)}{r^k},
        \qquad
        M(r)=\sup_{|z|=r}|y(z^{1/\mu})|.
\end{equation}
\Source{ThePadeMuntz\ldots Thm. \texttt{thm:radius}; RH-view Thm. ``Strictly positive Puiseux radius and Cauchy bound''; frmp.tex uses this as the analytic input behind the cgf lattice.}
\end{theorem}

\begin{proof}[Proof sketch]
On $B_1=\{y\in C[0,T_0]:\|y\|_\infty\le1\}$ define
\[
        \mathcal{V}(y)=I^{\mu}(c_0+c_1y+c_2y^2).
\]
The choice \eqref{eq:T0} makes $\mathcal{V}$ map $B_1$ into itself and makes the Lipschitz constant strictly below one.
Banach's fixed point theorem gives a unique solution.
Analytic dependence on $z$ follows by the same coefficient recurrence and the Cauchy majorant.
\end{proof}

\begin{theorem}[Generic finite radius]\label{thm:generic-finite}
Assume $c_2\ne0$ and $(c_0,c_1)\ne(0,0)$.
For generic values of $(c_0,c_1,c_2)$,
the radius $R$ in the $z$ plane is finite.
It equals the distance from the origin to the zero set of the entire Cole--Hopf denominator $w$ constructed in Theorem~\ref{thm:cole-hopf-explicit}.
\Source{ThePadeMuntz\ldots Thm. \texttt{thm:radius}; RH-view Thm. ``Generic finiteness of the radius''. The Cole--Hopf proof is the sharpest explanation.}
\end{theorem}

\begin{proof}[Proof sketch]
Theorem~\ref{thm:cole-hopf-explicit} writes $w$ as a nonzero linear combination of Mittag--Leffler functions.
Unless the combination reduces to a nonvanishing constant,
standard zero-distribution facts for finite-order entire functions give zeros at finite distance for a generic parameter set.
The solution is $g=-\Dc w/(c_2w)$,
so the first zero of $w$ determines the first pole of $g$ and hence the radius.
\end{proof}

\begin{remark}[Discrepancy resolved]
Some intermediate drafts suggest global convergence of the Puiseux series itself.
The correct statement is local convergence with positive radius,
followed by global meromorphic continuation and Padé convergence away from the Stahl compact or true poles.
The local series and the global rational approximants are different mathematical objects.
\end{remark}

\section{Cole--Hopf linearisation and meromorphy}

\begin{theorem}[Cole--Hopf linearisation]\label{thm:cole-hopf}
Let $g(z)=y(z^{1/\mu})$ solve
\begin{equation}\label{eq:riccati-z}
        \Dc g=c_0+c_1g+c_2g^2,
        \qquad g(0)=0.
\end{equation}
Assume $c_2\ne0$.
Define $w$ by
\begin{equation}\label{eq:w-def}
        \Dc w=-c_2gw,
        \qquad w(0)=1.
\end{equation}
Then
\begin{equation}\label{eq:g-cole}
        g=-\frac{\Dc w}{c_2w},
\end{equation}
and $w$ satisfies
\begin{equation}\label{eq:w-linear}
        \Dc^2w-c_1\Dc w+c_0c_2w=0,
        \qquad w(0)=1,
        \qquad \Dc w(0)=0.
\end{equation}
\Source{ThePadeMuntz\ldots Thm. \texttt{thm:colehopf}; RH-view Prop. ``Meromorphy of the transformed solution''; TheFractionalRiccatiMuntzPadeTheorem.tm Lem. \texttt{lem:meromorphic}.}
\end{theorem}

\begin{proof}
Equation \eqref{eq:w-def} gives \eqref{eq:g-cole} where $w\ne0$.
Apply $\Dc$ to $c_2gw=-\Dc w$.
Use Lemma~\ref{lem:muntz-leibniz} and \eqref{eq:riccati-z}:
\[
        c_2(\Dc g)w+c_2g(\Dc w)=-\Dc^2w.
\]
Since $\Dc w=-c_2gw$,
the quadratic terms cancel and yield \eqref{eq:w-linear}.
The initial conditions follow from $g(0)=0$ and $w(0)=1$.
\end{proof}

\begin{theorem}[Explicit entire Cole--Hopf denominator]\label{thm:cole-hopf-explicit}
Let
\begin{equation}\label{eq:lambdas}
        \lambda_{\pm}=\frac12\left(c_1\pm\sqrt{c_1^2-4c_0c_2}\right).
\end{equation}
If $\lambda_+\ne\lambda_-$, then
\begin{equation}\label{eq:w-explicit}
        w(z)=A\ML{\mu}{1}(\lambda_+z)+B\ML{\mu}{1}(\lambda_-z),
        \qquad
        A=\frac{-\lambda_-}{\lambda_+-\lambda_-},
        \quad
        B=\frac{\lambda_+}{\lambda_+-\lambda_-}.
\end{equation}
If $\lambda_+=\lambda_-=:\lambda$, then
\begin{equation}\label{eq:w-degenerate}
        w(z)=\ML{\mu}{1}(\lambda z)-\lambda z\ML{\mu}{1+\mu}(\lambda z).
\end{equation}
In both cases $w$ is entire of order $1/\mu$ in $z$.
\Source{ThePadeMuntz\ldots Thm. \texttt{thm:w\_entire}; RH-view rough-Heston meromorphy subsection.}
\end{theorem}

\begin{proof}
The Mittag--Leffler function
\[
        \ML{\mu}{1}(\lambda z)=\sum_{k\ge0}\frac{(\lambda z)^k}{\Gamma(k\mu+1)}
\]
satisfies $\Dc\ML{\mu}{1}(\lambda z)=\lambda\ML{\mu}{1}(\lambda z)$.
The nondegenerate formula is the unique linear combination with $w(0)=1$ and $\Dc w(0)=0$.
For the repeated root,
\[
        \Dc\{z\ML{\mu}{1+\mu}(\lambda z)\}=\ML{\mu}{1}(\lambda z),
\]
which verifies \eqref{eq:w-degenerate}.
The order statement is the standard Mittag--Leffler order computation.
\end{proof}

\begin{theorem}[Meromorphic continuation]\label{thm:meromorphy}
The transformed solution $g(z)=y(z^{1/\mu})$ extends meromorphically to the entire $z$ plane.
Its poles are the zeros of $w$.
They have no finite accumulation point.
A pole is simple when the zero of $w$ is simple;
in the constant-coefficient Riccati setting the zeros that generate poles are simple under the nondegeneracy condition
$w(z_0)=0\Rightarrow \Dc w(z_0)\ne0$.
\Source{ThePadeMuntz\ldots Thm. \texttt{thm:meromorphic}; RH-view Prop. \texttt{prop:meromorphy}. The simplicity clause is stated with the explicit nondegeneracy condition to avoid overclaiming at exceptional multiple zeros.}
\end{theorem}

\begin{proof}
The quotient formula \eqref{eq:g-cole} writes $g$ as a quotient of entire functions.
Since $w(0)=1$,
$w$ is not identically zero.
The identity theorem gives a discrete zero set with no finite accumulation point.
At a simple zero of $w$ the quotient has a simple pole unless the numerator vanishes there as well.
The displayed nondegeneracy prevents cancellation.
\end{proof}

\begin{remark}[Discrepancy resolved]
The clean method paper states that every pole is simple.
The quotient proof directly proves meromorphy and discreteness.
Simplicity follows for the generic Riccati parameters used by the Padé theory and for zeros with $\Dc w\ne0$.
At exceptional multiple zeros one must check cancellation or multiplicity from the entire functions themselves.
The canonical theorem therefore states the verified generic condition.
\end{remark}

\section{Padé approximation and Hankel systems}

\begin{definition}[Diagonal Padé approximant in the Müntz variable]\label{def:pade}
Let
\[
        g(z)=\sum_{k\ge0}m_kz^k
\]
be analytic at the origin.
The diagonal approximant $R_M=P_M/Q_M$ is characterized by
\begin{equation}\label{eq:pade-condition}
        Q_M(0)=1,
        \qquad
        \deg P_M\le M,
        \qquad
        \deg Q_M\le M,
        \qquad
        Q_Mg-P_M=O(z^{2M+1}).
\end{equation}
\Source{frmp.tex Padé section; ThePadeMuntz\ldots Lem. \texttt{lem:hankel}; RH-view Def. ``Diagonal Padé approximant and its domain''.}
\end{definition}

\begin{lemma}[Hankel denominator equations]\label{lem:hankel-denominator}
Writing $Q_M(z)=1+q_1z+\cdots+q_Mz^M$,
the Padé condition is equivalent to
\begin{equation}\label{eq:hankel-system}
        \sum_{j=1}^{M}m_{M+i-j}q_j=-m_{M+i},
        \qquad i=1,\ldots,M.
\end{equation}
The coefficient matrix is the Hankel matrix built from $m_1,\ldots,m_{2M-1}$ after the usual reversal of indices.
\Source{ThePadeMuntz\ldots Lem. \texttt{lem:hankel}; RH-view Lem. \texttt{lem:hankel-system}; FractionalRiccatiSolutionMethodology.tex Thm. \texttt{thm:hankelequiv}.}
\end{lemma}

\begin{proof}
Equate to zero the coefficients of $z^{M+1},\ldots,z^{2M}$ in $Q_Mg$.
The coefficient of $z^{M+i}$ is $m_{M+i}+\sum_{j=1}^Mq_jm_{M+i-j}$.
This gives \eqref{eq:hankel-system}.
\end{proof}

\begin{theorem}[Orthogonal-polynomial denominator]\label{thm:op-denominator}
Define the moment functional $\Lc$ by $\Lc[x^k]=m_k$.
Assume $\tau_n=\det(m_{i+j})_{0\le i,j<n}\ne0$ for $1\le n\le M$.
Let $p_M$ be the monic polynomial of degree $M$ satisfying
\begin{equation}\label{eq:op-orth}
        \Lc[x^kp_M(x)]=0,
        \qquad k=0,\ldots,M-1.
\end{equation}
Then the Padé denominator is
\begin{equation}\label{eq:den-reversal}
        Q_M(z)=z^Mp_M(1/z).
\end{equation}
The numerator is obtained by keeping the terms of degree at most $M$ in $Q_Mg$.
\Source{frmp.tex Thm. \texttt{thm:match}; FractionalRiccatiSolutionMethodology.tex Thms. \texttt{thm:padedenom}, \texttt{thm:padenum}, \texttt{thm:padeid}; TheFractionalRiccatiMuntzPadeTheorem.tm Prop. \texttt{prop:OPdenominator}; Gragg theorem cited through \cite{graggThePadeTableAndItsRelationToCertainAlgorithms}.}
\end{theorem}

\begin{proof}
Write $p_M(x)=\sum_{i=0}^M\pi_ix^i$ with $\pi_M=1$.
Then $Q_M(z)=\sum_{i=0}^M\pi_i z^{M-i}$.
For $0\le\ell<M$,
the coefficient of $z^{M+\ell}$ in $Q_Mg$ is
\[
        \sum_{i=0}^M\pi_i m_{i+\ell}=\Lc[p_Mx^\ell]=0.
\]
Thus all required middle coefficients vanish.
The remaining lower-degree part is the numerator.
\end{proof}

\begin{theorem}[Remainder structure]\label{thm:remainder}
For $g$ analytic at the origin and regular Padé order $M$,
\begin{equation}\label{eq:remainder}
        g(z)-R_M(z)=\frac{\Pi_{M+1}(z)}{Q_M(z)^2}z^{2M+1},
\end{equation}
where $\Pi_{M+1}$ is a polynomial of degree at most $M+1$ in the local algebraic case,
and analytic near the origin in the meromorphic case.
\Source{RH-view Thm. \texttt{thm:remainder-structure}; frmp.tex Thm. \texttt{thm:match}. The polynomial degree statement belongs to the algebraic local model; the meromorphic Riccati quotient gives the analytic version.}
\end{theorem}

\begin{proof}[Proof sketch]
The Padé identity gives a zero of order $2M+1$ at the origin for $Q_Mg-P_M$.
Divide by $Q_M$.
The second factor $Q_M$ in the denominator enters by the standard determinant representation of diagonal Padé remainders.
When the local continuation is algebraic,
clearing the algebraic relation yields the stated polynomial degree.
\end{proof}

\begin{theorem}[Global Padé--Müntz representation on positive time]\label{thm:global-positive}
Assume the Riccati solution has no pole on $[0,T^{\mu}]$ and that the diagonal Padé table satisfies the Gončar condition on that interval.
Then
\begin{equation}\label{eq:global-positive}
        y(t)=\lim_{M\to\infty}R_M(t^{\mu})
\end{equation}
uniformly for $0\le t\le T$.
\Source{ThePadeMuntz\ldots Thm. \texttt{thm:pade\_riccati}; RH-view Lem. \texttt{lem:goncar-interval} and Thm. \texttt{thm:global}. RH-view gives the sharpest hypotheses.}
\end{theorem}

\begin{proof}[Proof sketch]
Meromorphy follows from Theorem~\ref{thm:meromorphy}.
No true pole lies on the interval by hypothesis.
The Gončar condition controls spurious pole-zero pairs on the compact interval.
Gončar's theorem then promotes convergence in capacity to uniform convergence on the interval.
\end{proof}

\section{Orthogonal polynomials, Jacobi recurrences, and Fox--Wright structure}

\begin{theorem}[Determinant form and norms]\label{thm:det-op}
Assume $\tau_n\ne0$.
The monic polynomial $p_n$ satisfying \eqref{eq:op-orth} is
\begin{equation}\label{eq:op-det}
        p_n(x)=\frac1{\tau_n}
        \det\begin{pmatrix}
        m_0&m_1&\cdots&m_n\\
        m_1&m_2&\cdots&m_{n+1}\\
        \vdots&\vdots&&\vdots\\
        m_{n-1}&m_n&\cdots&m_{2n-1}\\
        1&x&\cdots&x^n
        \end{pmatrix},
\end{equation}
and
\begin{equation}\label{eq:norm-h}
        h_n=\Lc[x^np_n]=\frac{\tau_{n+1}}{\tau_n}\ne0.
\end{equation}
\Source{frmp.tex Lem. \texttt{lem:exist}; FractionalRiccatiSolutionMethodology.tex Defs. \texttt{def:Lop}, \texttt{def:Hankel}, \texttt{def:OPS}.}
\end{theorem}

\begin{proof}
Expanding along the last row shows monicity.
Applying $\Lc[x^k\cdot]$ with $k<n$ duplicates a row,
so orthogonality follows.
For $k=n$ the determinant is $\tau_{n+1}$.
Uniqueness is the nonsingularity of the Hankel system.
\end{proof}

\begin{theorem}[Three-term recurrence]\label{thm:three-term}
The monic orthogonal polynomials satisfy
\begin{equation}\label{eq:three-term-canonical}
        p_{n+1}(x)=(x-\alpha_n)p_n(x)-\beta_np_{n-1}(x),
        \qquad p_{-1}=0,
        \quad p_0=1,
\end{equation}
where
\begin{equation}\label{eq:alpha-beta}
        \alpha_n=\frac{\Lc[xp_n^2]}{h_n},
        \qquad
        \beta_n=\frac{h_n}{h_{n-1}}\quad(n\ge1),
        \qquad
        \beta_0=0.
\end{equation}
\Source{frmp.tex Thm. \texttt{thm:3term}; FractionalRiccatiSolutionMethodology.tex Cor. \texttt{cor:ABC}; TodaLattice.tex Jacobi section.}
\end{theorem}

\begin{proof}
Expand $xp_n$ in the orthogonal basis.
Pairing with $p_m$ gives zero for $m\le n-2$.
The coefficients of $p_n$ and $p_{n-1}$ are obtained by pairing with $p_n$ and $p_{n-1}$ respectively.
\end{proof}

\begin{theorem}[Chebyshev--Wheeler sigma table]\label{thm:cheb-wheeler}
Let
\begin{equation}\label{eq:sigma-def}
        \sigma_{k,\ell}=\Lc[p_kx^\ell],
        \qquad k\ge-1,
        \quad \ell\ge0.
\end{equation}
Set $\sigma_{-1,\ell}=0$ and $\sigma_{0,\ell}=m_\ell$.
Then
\begin{equation}\label{eq:sigma-rec}
        \sigma_{k+1,\ell}=\sigma_{k,\ell+1}-\alpha_k\sigma_{k,\ell}-\beta_k\sigma_{k-1,\ell}.
\end{equation}
The recurrence coefficients are recovered by
\begin{equation}\label{eq:sigma-extract}
        \alpha_0=\frac{m_1}{m_0},
        \qquad
        h_0=m_0,
\end{equation}
\begin{equation}\label{eq:sigma-extract-k}
        \alpha_k=\frac{\sigma_{k,k+1}}{\sigma_{k,k}}
        -\frac{\sigma_{k-1,k}}{\sigma_{k-1,k-1}},
        \qquad
        \beta_k=\frac{\sigma_{k,k}}{\sigma_{k-1,k-1}},
        \qquad
        h_k=\sigma_{k,k}.
\end{equation}
The construction costs $O(M^2)$ ring operations for orders up to $M$ and forms no Hankel inverse.
\Source{frmp.tex Thm. \texttt{thm:cheb}; FractionalRiccatiSolutionMethodology.tex Thm. \texttt{thm:CW}; frmp\_orthopoly.tex computational section; Wheeler algorithm \cite{wheelerModifiedMomentsGaussianQuadratures}.}
\end{theorem}

\begin{proof}
Multiply \eqref{eq:three-term-canonical} by $x^\ell$ and apply $\Lc$.
This gives \eqref{eq:sigma-rec}.
Orthogonality implies $h_k=\sigma_{k,k}$.
Solving the two displayed equations at the leading admissible powers gives \eqref{eq:sigma-extract-k}.
\end{proof}

\begin{definition}[Fox--Wright moment functional]\label{def:fox-wright-functional}
In the fractional Riccati setting,
\begin{equation}\label{eq:fw-moments}
        \Lc[x^k]=a_{k+1}(u)
\end{equation}
or,
for the cgf functional,
\begin{equation}\label{eq:fw-moments-d}
        \Lc[x^k]=d_k(v).
\end{equation}
The Fox--Wright deformation enters through the gamma-weight sequence
\begin{equation}\label{eq:psi-mu}
        \Psi_{\mu}(z)=\sum_{k\ge0}\frac{\Gamma(k\mu+1)}{\Gamma(k\mu+\mu+1)}z^k
        ={}_1\Psi_1\!\left(\begin{matrix}(1,\mu)\\(1+\mu,\mu)\end{matrix};z\right).
\end{equation}
\Source{frmp\_orthopoly.tm Def. ``Fox--Wright moment functional''; frmp\_orthopoly.tex Def. ``Fox--Wright moment functional'' and Fox--Wright deformation subsection.}
\end{definition}

\begin{theorem}[Classification of the fractional Riccati orthogonal family]\label{thm:fw-classification}
For generic complex rough-Heston parameter $u$,
the polynomials $p_n(\cdot,u)$ are quasi-definite non-Hermitian orthogonal polynomials.
They are not classical Jacobi,
Laguerre,
or Hermite polynomials in general.
Their Stahl-limit orthogonality lives on the minimal-capacity compact $K_A(u)$,
with jump weight obtained from the Fox--Wright-deformed discriminant
\begin{equation}\label{eq:fw-discriminant}
        \Delta_F(z,u)=(1-zQ(u))^2-4zP(u)R(u)\Psi_{\mu}(z).
\end{equation}
Across the compact,
the jump has the local form
\begin{equation}\label{eq:fw-jump}
        \rho_F(z,u)=\frac{\sqrt{\Delta_F^+(z,u)}-\sqrt{\Delta_F^-(z,u)}}{2zR(u)}
	times \text{a meromorphic factor fixed by the Riccati branch}.
\end{equation}
The zeros of $p_n$ equidistribute according to the equilibrium measure on $K_A(u)$.
\Source{frmp\_orthopoly.tm classification section; frmp\_orthopoly.tex Sections ``Classification'' and ``Asymptotic Theory''. This result appears only in the orthopoly article and is incorporated here as the canonical classification.}
\end{theorem}

\begin{proof}[Proof sketch]
The moment sequence is generated by the Riccati recurrence with gamma ratios.
Those ratios assemble into \eqref{eq:psi-mu}.
The algebraic Riccati relation for the auxiliary branch gives \eqref{eq:fw-discriminant}.
Taking boundary values across the Stahl compact gives the jump formula.
Stahl--Nuttall theory then identifies the zero-counting measures with the equilibrium measure.
\end{proof}

\begin{proposition}[Mittag--Leffler and Laguerre corner]\label{prop:ml-laguerre}
When the quadratic coefficient $R(u)$ vanishes,
the Riccati recurrence becomes linear and the associated moment functional reduces to a Mittag--Leffler moment sequence.
The resulting orthogonal system is a Laguerre-type degeneration of the Fox--Wright non-Hermitian family.
\Source{frmp\_orthopoly.tm Section ``$R=0$ Oracle''; frmp\_orthopoly.tex Section ``$R=0$ Oracle: Mittag-Leffler Moments and Laguerre Structure''.}
\end{proposition}

\begin{proof}[Proof sketch]
With $R=0$ the convolution term disappears.
The moments are produced by repeated multiplication by $Q$ and gamma ratios.
These are precisely Mittag--Leffler moments,
whose inverse Laplace representation gives the Laguerre-type limiting structure.
\end{proof}

\section{Riemann--Hilbert, Stahl, and global convergence theory}

\begin{definition}[Cauchy transform and FIK problem]\label{def:fik}
Let $\Sigma$ be a finite union of oriented arcs and $\omega$ an integrable complex weight.
The Cauchy transform is
\begin{equation}\label{eq:cauchy-transform}
        f(z)=\int_{\Sigma}\frac{\omega(s)\dd s}{z-s}.
\end{equation}
The Fokas--Its--Kitaev Riemann--Hilbert problem seeks $Y:\C\setminus\Sigma\to GL_2(\C)$ with jump
\begin{equation}\label{eq:fik-jump}
        Y_+(s)=Y_-(s)
        \begin{pmatrix}1&2\pi i\omega(s)\\0&1\end{pmatrix},
\end{equation}
and normalization
\begin{equation}\label{eq:fik-normal}
        Y(z)=\left(I+O(z^{-1})\right)\diag(z^n,z^{-n}).
\end{equation}
Then $Y_{11}$ is the degree-$n$ denominator and $Y_{12}/Y_{11}$ gives the Padé approximant.
\Source{RH-view Defs. \texttt{def:cauchy}, \texttt{def:pade}, Thm. \texttt{thm:fik}, Cor. \texttt{cor:orthogonality}.}
\end{definition}

\begin{theorem}[FIK representation after selecting the contour]\label{thm:fik}
Once the Stahl compact $K_g$ for a branch of $g$ has been selected,
the diagonal Padé denominator is the first entry of the solution to the FIK problem on $K_g$ with weight equal to the jump of $g$ across $K_g$.
The orthogonality relations are
\begin{equation}\label{eq:fik-orth}
        \int_{K_g}s^kQ_n(s)\omega_g(s)\dd s=0,
        \qquad k=0,\ldots,n-1.
\end{equation}
\Source{RH-view Thm. \texttt{thm:fik} and Remark \texttt{rem:fixed-contour}; frmp\_orthopoly.tex non-Hermitian orthogonality section.}
\end{theorem}

\begin{proof}[Proof sketch]
The FIK theorem is standard for a fixed contour.
The present functions are initially specified by local series,
so the contour is first determined by Stahl's extremal problem.
On that contour the jump of $g$ supplies the weight,
and the FIK solution is the usual orthogonal-polynomial matrix.
\end{proof}

\begin{definition}[Stahl compact]\label{def:stahl}
A function analytic at infinity belongs to the Stahl class if it has multivalued meromorphic continuation outside a compact set of logarithmic capacity zero.
An admissible compact is a compact set outside which one branch is single-valued and meromorphic.
The Stahl compact $K_g$ is the admissible compact of minimal capacity.
\Source{RH-view Defs. \texttt{def:stahl-class}, \texttt{def:admissible}; Stahl references \cite{stahlExtremalDomainsAssociatedWithAnAnalyticFunctionI,stahlExtremalDomainsAssociatedWithAnAnalyticFunctionII,stahlConvergenceOfPadeApproximantsBranchPoints}.}
\end{definition}

\begin{theorem}[Stahl compact and $S$-property]\label{thm:stahl-compact}
For $g$ in the Stahl class,
there is a unique compact $K_g$ of minimal logarithmic capacity.
If $D_g=\widehat{\C}\setminus K_g$ and $G(z)$ is the Green function of $D_g$ with pole at infinity,
then the arcs of $K_g$ satisfy
\begin{equation}\label{eq:s-property-canonical}
        \frac{\partial G}{\partial n_+}
        =
        \frac{\partial G}{\partial n_-}
\end{equation}
away from endpoints.
\Source{RH-view Thms. \texttt{thm:stahl-geometric}, \texttt{thm:s-property}; frmp\_orthopoly.tex Section ``The $S$-Property and Stahl's Compact''.}
\end{theorem}

\begin{proof}[Proof sketch]
The compact is obtained by minimizing logarithmic capacity over admissible cuts.
The equality of normal derivatives is the Euler--Lagrange equation for that variational problem.
Potential theory supplies uniqueness.
\end{proof}

\begin{theorem}[Stahl convergence and rate]\label{thm:stahl-convergence}
Let $g$ be in the Stahl class with compact $K_g$ and pole set $\mathcal{P}_g$ in $D_g$.
The diagonal Padé approximants converge to $g$ in capacity on compact subsets of $D_g\setminus\mathcal{P}_g$.
At regular points,
\begin{equation}\label{eq:stahl-rate}
        \limsup_{M\to\infty}|g(z)-R_M(z)|^{1/M}\le e^{-G(z)}.
\end{equation}
For the cgf,
this gives convergence of $\Phi_M(v,T)$ whenever $T^{\mu}\in D_g$ and is not a pole.
\Source{frmp.tex Thm. \texttt{thm:stahl}; RH-view Thms. \texttt{thm:stahl-convergence}, \texttt{thm:deift-zhou}, Cor. \texttt{cor:exponential-rate}. RH-view supplies the sharp potential-theoretic hypotheses.}
\end{theorem}

\begin{proof}[Proof sketch]
Stahl's theorem gives convergence in capacity for diagonal Padé approximants of functions with finite-capacity branch structure.
The Green-function rate is the standard consequence of the same theorem,
refined by nonlinear steepest descent when the jump is analytic and endpoints are regular.
\end{proof}

\begin{theorem}[de Montessus, Nuttall--Pommerenke, and Gončar synthesis]\label{thm:local-global-synthesis}
In a disk containing only finitely many poles of $g$ and no branch cut,
de Montessus de Ballore gives locally uniform convergence of row sequences with fixed denominator degree to the meromorphic function.
For diagonal sequences of a multivalued function,
Nuttall--Pommerenke permits bounded pole-zero pairs away from the limiting compact.
If the number of approximant poles in a compact set is uniformly bounded,
Gončar's condition yields uniform convergence there.
\Source{ThePadeMuntz\ldots Thms. \texttt{thm:dMdB}, \texttt{thm:NP}; RH-view Def. \texttt{def:goncar}, Thm. \texttt{thm:goncar}, Lem. \texttt{lem:goncar-interval}; frmp.tex convergence section.}
\end{theorem}

\begin{proof}[Proof sketch]
The three named theorems apply at different geometric scales.
de Montessus controls finitely many stable poles.
Nuttall--Pommerenke controls exceptional local pole-zero pairs for diagonal Padé sequences.
Gončar turns bounded exceptional counts into uniform convergence on compact sets disjoint from true singularities.
\end{proof}

\begin{theorem}[A posteriori Padé certificate]\label{thm:apost}
Let $R_M$ be consecutive diagonal approximants at a fixed $z$,
and let
\begin{equation}\label{eq:deltaM}
        \Delta_M(z)=R_M(z)-R_{M-1}(z).
\end{equation}
Assume $|\Delta_M|<|\Delta_{M-1}|$ and the later differences decrease with a common ratio no larger than
\begin{equation}\label{eq:rho-star}
        \rho_* =\frac{|\Delta_M|}{|\Delta_{M-1}|}<1.
\end{equation}
Then
\begin{equation}\label{eq:apost-canonical}
        |g(z)-R_M(z)|
        \le
        \frac{|\Delta_M(z)|^2}{|\Delta_{M-1}(z)|-|\Delta_M(z)|}.
\end{equation}
\Source{frmp.tex Thm. \texttt{thm:apost}; ThePadeMuntz\ldots Thm. \texttt{thm:error}; RH-view Thm. \texttt{thm:error-bound}. RH-view's observable-ratio form is the canonical one.}
\end{theorem}

\begin{proof}
The difference telescoping identity gives
$g-R_M=\sum_{j\ge1}\Delta_{M+j}$.
The geometric assumption bounds this by
$\rho_*|\Delta_M|/(1-\rho_*)$,
which is exactly \eqref{eq:apost-canonical}.
\end{proof}

\section{Toda lattice, Jacobi operator, and spectral measures}

\begin{definition}[Jacobi matrix and finite sections]\label{def:jacobi}
From \eqref{eq:three-term-canonical} form
\begin{equation}\label{eq:jacobi}
        J=\begin{pmatrix}
        \alpha_0&1&0&\cdots\\
        \beta_1&\alpha_1&1&\ddots\\
        0&\beta_2&\alpha_2&\ddots\\
        \vdots&\ddots&\ddots&\ddots
        \end{pmatrix}.
\end{equation}
The $M\times M$ leading finite section is denoted $J_M$.
In the positive self-adjoint case one instead symmetrizes by using off-diagonal entries $\sqrt{\beta_n}$.
The associated spectral measure is denoted $\nu_g$.
\Source{RH-view Def. \texttt{def:jacobi}; FractionalRiccatiSolutionMethodology.tex Def. \texttt{def:Jacobi}; TodaLattice.tex Jacobi-operator sections.}
\end{definition}

\begin{theorem}[Padé poles as Jacobi eigenvalues]\label{thm:poles-eigenvalues}
The zeros of $Q_M(z)=z^Mp_M(1/z)$ are reciprocals of the eigenvalues of $J_M$ in the $x$ variable.
Equivalently,
the poles of $R_M(z)$ are the $z$-plane points associated with the spectrum of the finite Jacobi section.
In the self-adjoint Cauchy-transform case,
\begin{equation}\label{eq:resolvent-rep}
        R_M(z)=\langle e_0,(J_M-zI)^{-1}e_0\rangle
        =\sum_{j=1}^M\frac{\omega_j^{(M)}}{\lambda_j^{(M)}-z},
\end{equation}
where $\omega_j^{(M)}=|\langle e_0,\psi_j^{(M)}\rangle|^2$.
\Source{RH-view Thm. \texttt{thm:resolvent-representation}, Cor. \texttt{cor:poles-eigenvalues}; FractionalRiccatiSolutionMethodology.tex Thm. \texttt{thm:spectrum}, Cor. \texttt{cor:poles}.}
\end{theorem}

\begin{proof}[Proof sketch]
The three-term recurrence is the characteristic-polynomial recurrence for the finite Jacobi section.
Thus the zeros of $p_M$ are eigenvalues of $J_M$.
The reciprocal map converts $p_M$ zeros into $Q_M$ zeros.
For a positive measure,
the spectral theorem gives \eqref{eq:resolvent-rep}.
\end{proof}

\begin{theorem}[Toda flow of recurrence coefficients]\label{thm:toda}
Let the moment functional evolve by
\begin{equation}\label{eq:moment-flow}
        \dd\nu_s(x)=e^{sx}\dd\nu_0(x)
\end{equation}
in the self-adjoint case.
Then the recurrence coefficients satisfy the Toda equations
\begin{equation}\label{eq:toda-canonical}
        \dot{\alpha}_n=\beta_{n+1}-\beta_n,
        \qquad
        \dot{\beta}_n=\beta_n(\alpha_n-\alpha_{n-1}).
\end{equation}
The spectrum of the infinite Jacobi operator is invariant along the flow.
\Source{RH-view Thm. \texttt{thm:toda}; TodaLattice.tex Lax-pair, Jacobi-operator, and spectral-map sections.}
\end{theorem}

\begin{proof}[Proof sketch]
Multiplication of the measure by $e^{sx}$ differentiates moments by shifting the index.
The compatibility equations for the orthogonal-polynomial recurrence and the $s$ evolution give \eqref{eq:toda-canonical}.
Equivalently,
$\dot J=[B,J]$ for the skew part $B$ of $J$.
A Lax flow preserves spectrum.
\end{proof}

\begin{theorem}[Spectral reconstruction certificate]\label{thm:spectral-certificate}
The Hankel determinant route,
the Chebyshev--Wheeler sigma table,
and the Jacobi finite-section route compute the same denominator whenever $\tau_M\ne0$.
The a posteriori certificate of Theorem~\ref{thm:apost} can be interpreted as a spectral-gap certificate for the change in resolvents of successive finite sections.
\Source{RH-view Thm. \texttt{thm:hankel-jacobi}, Cor. \texttt{cor:spectral-gap}, Thm. \texttt{thm:compiled-spectral}; FractionalRiccatiSolutionMethodology.tex Thms. \texttt{thm:hankelequiv}, \texttt{thm:algcorr}.}
\end{theorem}

\begin{proof}[Proof sketch]
All three constructions are equivalent to the same orthogonality relations.
Successive finite sections change the resolvent by a rank-one boundary update.
The difference of the scalar resolvents is precisely the Padé difference used in Theorem~\ref{thm:apost}.
\end{proof}

\section{Rough Heston pricing}

\begin{definition}[Rough-Heston transform data]\label{def:heston-data}
Let $S_T=S_0e^{X_T+rT}$.
The rough-Heston characteristic function has the affine form
\begin{equation}\label{eq:char-affine}
        \phi(v,T)=\exp\Phi(v,T),
\end{equation}
where $\Phi$ is assembled from the fractional Riccati solution with coefficients
\begin{equation}\label{eq:heston-symbols}
        c_0=P(v),
        \qquad
        c_1=Q(v),
        \qquad
        c_2=R(v).
\end{equation}
The Padé-resummed cgf is
\begin{equation}\label{eq:heston-pade}
        \Phi_M(v,T)=T^{\mu}\frac{P_M(T^{\mu};v)}{Q_M(T^{\mu};v)}.
\end{equation}
\Source{frmp.tex Sections \texttt{sec:standing}, \texttt{sec:cgf}, \texttt{sec:pade}; FractionalRiccatiSolutionMethodology.tex rough-Heston section; TheFractionalRiccatiMuntzPadeTheorem.tm application section.}
\end{definition}

\begin{theorem}[Lewis contour representation]\label{thm:lewis}
Assume the moment condition $p^*>1$ and choose $c\in(1,p^*)$.
For the call payoff,
\begin{equation}\label{eq:lewis-call}
        C_M(K,T)=\frac{Ke^{-rT}}{2\pi i}
        \int_{c-i\infty}^{c+i\infty}
        \frac{\exp(\Phi_M(-iw,T))e^{-w\tilde k}}{w(w-1)}\dd w,
\end{equation}
where $\tilde k=\log(K/S_0)-rT$ under the stated normalization.
The put is obtained by put--call parity.
\Source{frmp.tex Lem. \texttt{lem:payoff}, Thm. \texttt{thm:lewis}; residue\_price.tex Thm. \texttt{thm:main}; single-series-pricing.md Prop. 2.1.}
\end{theorem}

\begin{proof}[Proof sketch]
The bilateral Laplace transform of the call payoff is $1/(w(w-1))$ on the strip $\Re w>1$.
Parseval inversion for bilateral Laplace transforms yields \eqref{eq:lewis-call}.
The moment condition places the contour inside the analyticity strip.
\end{proof}

\begin{lemma}[Schwinger--Gauss--erfc kernel]\label{lem:erfc-kernel}
Let
\begin{equation}\label{eq:g-gaussian}
        G(w)=\frac12\sigma_T^2w^2-\xi w+\chi
\end{equation}
with $\sigma_T^2>0$.
For $\Re w_*<c$,
\begin{equation}\label{eq:erfc-kernel}
        \mathcal{E}(w_*)
        :=\frac1{2\pi i}\int_{c-i\infty}^{c+i\infty}
        \frac{e^{G(w)}}{w-w_*}\dd w
        =\frac12e^{G(w_*)}\erfc\left(\frac{\xi-\sigma_T^2w_*}{\sigma_T\sqrt2}\right).
\end{equation}
For $\Re w_*>c$ one subtracts the crossed residue $e^{G(w_*)}$.
Moreover,
\begin{equation}\label{eq:hermite-atoms}
        \mathcal{A}_m(w_*)=\frac1{2\pi i}\int_{c-i\infty}^{c+i\infty}
        \frac{e^{G(w)}}{(w-w_*)^m}\dd w
        =\frac1{(m-1)!}\partial_{w_*}^{m-1}\mathcal{E}(w_*),
\end{equation}
and the derivatives are finite combinations of $\erfc$ and Hermite functions.
\Source{frmp.tex Lemmas \texttt{lem:E}, \texttt{lem:hermite}, \texttt{lem:dE}; single-series-pricing.md Lemmas 4.1--4.2.}
\end{lemma}

\begin{proof}
For $\Re w_*<c$ write
$(w-w_*)^{-1}=\int_0^{\infty}e^{-(w-w_*)s}\dd s$.
The inner contour integral is Gaussian after completing the square.
The remaining half-line Gaussian integral is the complementary error function.
Higher-order kernels follow by differentiating with respect to $w_*$.
\end{proof}

\begin{theorem}[Residue expansion from rational exponent]\label{thm:residue-price}
Suppose the rational part of $\Phi_M(-iw,T)-w\tilde k$ has simple poles $\eta_j$ with residues $r_j$,
and suppose the Gaussian coefficient satisfies $\sigma_T^2>0$.
Then the Lewis integral can be evaluated by expanding the pole exponentials and applying the kernel atoms \eqref{eq:hermite-atoms}.
When repeated poles occur,
Heaviside cover-up derivatives replace simple residues.
\Source{frmp.tex Lemmas \texttt{lem:rotate}, \texttt{lem:cauchy}, \texttt{lem:heav}, Thm. \texttt{thm:price}; residue\_price.tex Thm. \texttt{thm:main}.}
\end{theorem}

\begin{proof}[Proof sketch]
Separate the Gaussian quadratic part from the remaining rational exponent.
Each pole contribution is expanded as an exponential of a principal part.
The product of those absolutely convergent local expansions is integrated termwise under the Gaussian majorant.
Each term is a finite linear combination of \eqref{eq:hermite-atoms}.
\end{proof}

\begin{theorem}[One-index erfc--Hermite pricing series]\label{thm:single-series}
Choose $c\in(1,p^*)$ and $\kappa>0$.
Set
\begin{equation}\label{eq:cayley}
        \zeta(w)=\frac{w-c+\kappa}{w-c-\kappa},
        \qquad
        w(\zeta)=c-\kappa\frac{1+\zeta}{1-\zeta},
        \qquad q=c+\kappa.
\end{equation}
If all singularities of the non-Gaussian exponent lie in $\Re w>c$,
then their images under $\zeta$ lie outside the unit disk.
Writing
\begin{equation}\label{eq:q-series}
        \Psi(\zeta)=\exp\{\text{non-Gaussian rational exponent at }w(\zeta)\}
        =\sum_{n\ge0}q_n\zeta^n,
\end{equation}
the call price has the absolutely convergent single-index expansion
\begin{equation}\label{eq:single-price}
        C_M(K,T)=Ke^{-rT}\sum_{n\ge0}q_n\left(T_n^{(1)}-T_n^{(0)}\right),
\end{equation}
where
\begin{equation}\label{eq:Tn-delta}
        T_n^{(\delta)}=\frac1{2\pi i}\int_{c-i\infty}^{c+i\infty}
        \frac{e^{G(w)}\zeta(w)^n}{w-\delta}\dd w,
        \qquad \delta\in\{0,1\}.
\end{equation}
The coefficients satisfy
\begin{equation}\label{eq:exp-series-rec}
        q_0=e^{r_0},
        \qquad
        nq_n=\sum_{m=1}^n m r_m q_{n-m},
\end{equation}
where $r_m$ are the Taylor coefficients of the rational exponent in $\zeta$.
The kernel terms are generated from the three fixed nodes $0,1,q$ by the erfc--Hermite atoms.
\Source{single-series-pricing.md Sections 5--8; frmp.tex kernel lemmas and Black--Scholes corollary. This is the preferred pricing-function form because it avoids root enumeration.}
\end{theorem}

\begin{proof}[Proof sketch]
The Cayley map sends the integration line to $|\zeta|=1$ and the analytic half-plane to $|\zeta|<1$.
Since the singularities are outside,
$\Psi$ has radius greater than one and the Cauchy bound gives uniform convergence on the contour.
Expanding $\zeta(w)^n=(1+2\kappa/(w-q))^n$ reduces every term to finitely many atoms at $0,1,q$.
Absolute convergence permits exchanging sum and integral.
The recurrence \eqref{eq:exp-series-rec} is the coefficient identity $\Psi'=R'\Psi$.
\end{proof}

\begin{corollary}[Black--Scholes limit]\label{cor:black-scholes}
If the non-Gaussian rational part vanishes and
$\mu_T=\tfrac12\sigma_T^2$,
then only the $n=0$ term remains in \eqref{eq:single-price},
and
\begin{equation}\label{eq:bs}
        C=S_0N(d_1)-Ke^{-rT}N(d_2),
        \qquad
        d_{1,2}=\frac{\log(S_0/K)+rT}{\sigma_T}\pm\frac{\sigma_T}{2}.
\end{equation}
\Source{frmp.tex Cor. \texttt{cor:BS}; single-series-pricing.md Section 8.}
\end{corollary}

\begin{proof}
With the non-Gaussian part equal to zero,
$q_0=1$ and $q_n=0$ for $n\ge1$.
The price is the difference of the two erfc kernels at $0$ and $1$.
Using $N(x)=\tfrac12\erfc(-x/\sqrt2)$ gives \eqref{eq:bs}.
\end{proof}

\begin{remark}[Pricing discrepancy resolved]
The multi-index residue expansion in \texttt{frmp.tex} is valid under its pole-separation and Gaussian-decay assumptions.
The one-index Cayley expansion in \texttt{single-series-pricing.md} proves the same contour integral under a condition expressed by the analyticity radius of $\Psi$ in the Cayley disk.
The canonical implementation route is the one-index expansion because its coefficients are obtained by scalar recurrences and its error bound does not require listing the Padé poles.
\end{remark}

\section{Error bounds and certified numerics}

\begin{theorem}[Padé-to-cgf error propagation]\label{thm:cgf-error}
Let $\varepsilon_M(v,T)$ bound
$|\Phi(v,T)-\Phi_M(v,T)|$
on a pricing contour.
Then
\begin{equation}\label{eq:exp-error}
        |e^{\Phi(v,T)}-e^{\Phi_M(v,T)}|
        \le e^{\Re\Phi_M(v,T)}\left(e^{\varepsilon_M(v,T)}-1\right)
\end{equation}
whenever $|\Phi-\Phi_M|\le\varepsilon_M$.
If $\varepsilon_M\le1$,
this is bounded by $e^{\Re\Phi_M+1}\varepsilon_M$.
\Source{frmp.tex Thm. \texttt{thm:ed-conv}; RH-view a-posteriori section.}
\end{theorem}

\begin{proof}
Use
$e^a-e^b=e^b(e^{a-b}-1)$
and
$|e^\delta-1|\le e^{|\delta|}-1$.
\end{proof}

\begin{theorem}[Single-series remainder certificate]\label{thm:single-error}
In the setting of Theorem~\ref{thm:single-series},
let $1<R<R_0$ be inside the Cayley analyticity radius and set
\begin{equation}\label{eq:GN}
        G_\delta=\frac1{2\pi}\int_{\R}\frac{|e^{G(c+i\xi)}|}{|c-\delta+i\xi|}\dd \xi.
\end{equation}
If $|q_n|\le A R^{-n}$,
then the partial sum
\begin{equation}\label{eq:SN}
        S_N=Ke^{-rT}\sum_{n=0}^{N}q_n(T_n^{(1)}-T_n^{(0)})
\end{equation}
satisfies
\begin{equation}\label{eq:single-error}
        |C_M-S_N|
        \le
        Ke^{-rT}(G_0+G_1)A\frac{R^{-(N+1)}}{1-R^{-1}}.
\end{equation}
\Source{single-series-pricing.md Theorem 7.1 and Section 7.2.}
\end{theorem}

\begin{proof}
On the contour $|\zeta|=1$.
Therefore $|T_n^{(\delta)}|\le G_\delta$ uniformly in $n$.
The geometric coefficient bound gives \eqref{eq:single-error} by summing the tail of a geometric series.
\end{proof}

\begin{theorem}[End-to-end certified computation]\label{thm:end-to-end}
A certified computation of a rough-Heston price consists of the following verified inclusions:
\begin{enumerate}[label=(\roman*)]
\item ball enclosures for $a_k$ or $d_k$ from the Müntz--Tau recurrence;
\item ball enclosures for $\alpha_n,\beta_n,h_n$ from the Chebyshev--Wheeler sigma table;
\item a Padé enclosure using Theorem~\ref{thm:apost};
\item an exponential enclosure using Theorem~\ref{thm:cgf-error};
\item a pricing-series enclosure using Theorem~\ref{thm:single-error} or the residue expansion under its assumptions.
\end{enumerate}
The output is certified when the final ball radius is below the requested tolerance.
\Source{frmp.tex Cor. \texttt{cor:pipeline}; single-series-pricing.md implementation section; FractionalRiccatiSolutionMethodology.tex Thms. \texttt{thm:algcorr}, \texttt{thm:complexity}, Prop. \texttt{prop:verify}.}
\end{theorem}

\begin{proof}[Proof sketch]
Each step is an inclusion-preserving algebraic or analytic operation.
The recurrence steps are finite ball operations.
The Padé and pricing bounds are explicit inequalities.
Composing inclusions gives the final ball.
\end{proof}

\begin{remark}[Discrepancy resolved]
Several drafts describe Hankel solves as the main computational route.
The determinant formulas are mathematically canonical,
but the Chebyshev--Wheeler route is the stable algorithmic form because it computes the same recurrence data with quadratic arithmetic complexity and without Hankel inversion.
\end{remark}

\section{Implementation in arb4j}

\begin{definition}[Canonical data flow]\label{def:arb4j-flow}
The arb4j implementation should follow this data flow:
\begin{enumerate}[label=(\arabic*)]
\item build $P(v),Q(v),R(v)$ for the model;
\item compute $a_k(v)$ by Theorem~\ref{thm:tau};
\item derive $d_k(v)$ by Definition~\ref{def:cgf-coefficients};
\item build $\alpha_n,\beta_n,h_n$ by Theorem~\ref{thm:cheb-wheeler};
\item assemble $Q_M$ and $P_M$ by Theorem~\ref{thm:op-denominator};
\item evaluate $\Phi_M$ and the certified price by Theorems~\ref{thm:single-series} and \ref{thm:single-error}.
\end{enumerate}
\Source{frmp.tex pipeline; FractionalRiccatiSolutionMethodology.tex algorithm and complexity sections; single-series-pricing.md Section 10.}
\end{definition}

\begin{proposition}[Exact-arithmetic structure]\label{prop:exact-structure}
Before numerical evaluation at a fixed $v$,
the recurrence data live in the rational function field generated by the model coefficients and gamma constants.
After evaluation,
all computations are complex ball operations.
The last-argument-is-result style of arb4j arithmetic matches the recurrence form because every update has a designated output object.
\Source{arb4j implementation conventions; FractionalRiccatiSolutionMethodology.tex operation-count section.}
\end{proposition}

\begin{proof}[Proof sketch]
The recurrences use additions,
multiplications,
divisions by nonzero Hankel norms,
and gamma constants.
Those operations stay in the stated field symbolically and in ball arithmetic numerically.
\end{proof}

\begin{theorem}[Implementation correctness contract]\label{thm:implementation-contract}
An implementation conforms to the Fractional Riccati--Müntz--Padé theory if it satisfies:
\begin{enumerate}[label=(\roman*)]
\item the expression for the Riccati equation is literal mathematical source at the call site;
\item the coefficient arrays are represented by arb vector objects or polynomial objects, not by Java arrays of ball handles;
\item every native ball object has deterministic ownership and closing;
\item the Padé denominator is obtained by reciprocal orthogonal polynomials;
\item every returned price is accompanied by a ball radius from a theorem in Section 11.
\end{enumerate}
\Source{arb4j codebase conventions; frmp.tex certification theorem; single-series-pricing.md implementation directives.}
\end{theorem}

\begin{proof}[Proof sketch]
Items (i)--(iii) are representation invariants required by the expression compiler and native memory model.
Items (iv)--(v) are exactly the mathematical construction and certification chain proved above.
\end{proof}

\begin{remark}[Discrepancy resolved]
The source files alternate between denominator-first,
Jacobi-first,
and Riemann--Hilbert-first descriptions.
They are equivalent under quasi-definiteness.
The implementation should be denominator-first through Chebyshev--Wheeler,
while documentation can still expose the Jacobi and Riemann--Hilbert interpretations for analysis.
\end{remark}

\appendix
\section{Crosswalk from source files to canonical results}

\begin{longtable}{p{0.23\textwidth}p{0.42\textwidth}p{0.27\textwidth}}
\toprule
Source file & Results and role & Location in this exposition\\
\midrule
\endfirsthead
\toprule
Source file & Results and role & Location in this exposition\\
\midrule
\endhead
\texttt{frmp.tex} & Master pricing paper; fractional calculus; Müntz--Tau recurrence; rough-Heston cgf lattice; orthogonal polynomials; Chebyshev--Wheeler; Padé matching; Stahl rate; a posteriori bound; Lewis integral; erfc and Hermite kernels; multi-index residue expansion; Black--Scholes limit; end-to-end price convergence. & Defs. \ref{def:frac-ops}, \ref{def:cgf-coefficients}; Thms. \ref{thm:tau}, \ref{thm:cgf}, \ref{thm:det-op}, \ref{thm:three-term}, \ref{thm:cheb-wheeler}, \ref{thm:op-denominator}, \ref{thm:stahl-convergence}, \ref{thm:apost}, \ref{thm:lewis}, \ref{lem:erfc-kernel}, \ref{thm:residue-price}, \ref{cor:black-scholes}, \ref{thm:end-to-end}.\\
\texttt{ThePadeMuntzSpectralTauMethodForConstantCoeffecientFractionalRiccations.tex} & Cleanest scalar method paper; Riemann--Liouville and Caputo definitions; power rules; Müntz closure; series solution; positive radius; Cole--Hopf linearisation; explicit entire denominator; meromorphy; Hankel system; computable error; de Montessus and Nuttall--Pommerenke; global positive-axis representation. & Sections 1--6; Thms. \ref{thm:radius}, \ref{thm:generic-finite}, \ref{thm:cole-hopf}, \ref{thm:cole-hopf-explicit}, \ref{thm:meromorphy}, \ref{lem:hankel-denominator}, \ref{thm:global-positive}, \ref{thm:local-global-synthesis}.\\
\texttt{TheRiemannHilbertViewOfPadeApproximationAndGlobalAnalyticContinuation.tex} & Most complete general Padé theory; FIK Riemann--Hilbert formulation; Stahl compact; $S$-property; convergence in capacity; Froissart and Gončar criteria; Hankel tau functions; Jacobi operator; Toda flow; Fekete and zero condensation; bordered rank-one update; spectral reconstruction. & Sections 6--9; Defs. \ref{def:fik}, \ref{def:stahl}, \ref{def:jacobi}; Thms. \ref{thm:fik}, \ref{thm:stahl-compact}, \ref{thm:stahl-convergence}, \ref{thm:local-global-synthesis}, \ref{thm:toda}, \ref{thm:spectral-certificate}.\\
\texttt{frmp\_orthopoly.tm} & Fox--Wright orthogonal-polynomial family; classification as non-Hermitian OPs; Fox--Wright deformation; Mittag--Leffler and Laguerre corner; Müntz--Szász perspective; Chebyshev--Wheeler realization. & Def. \ref{def:fox-wright-functional}; Thm. \ref{thm:fw-classification}; Prop. \ref{prop:ml-laguerre}; Thm. \ref{thm:cheb-wheeler}.\\
\texttt{frmp\_orthopoly.tex} & LaTeX version of the Fox--Wright classification; explicit discriminant, jump weight, zero condensation, recurrence-coefficient asymptotics. & Def. \ref{def:fox-wright-functional}; Thm. \ref{thm:fw-classification}; Section 7.\\
\texttt{frmp\_pade\_muntz.tex} & Early closed-form pricing draft; supports the residue-price route and the interpretation of Padé poles and residues in the Fourier plane. & Thm. \ref{thm:residue-price}; Section 10 remarks.\\
\texttt{FractionalRiccatiSolutionMethodology.tex} & Intermediate algorithmic synthesis; polynomial degree growth; moment functional; Hankel determinants; signed quasi-definiteness; generating polynomials $S(n,z)$; Chebyshev--Wheeler polynomial recurrence; second-kind numerator; Jacobi spectrum; complexity; validation; correction notes. & Cor. \ref{cor:degree-growth}; Thms. \ref{thm:cheb-wheeler}, \ref{thm:op-denominator}, \ref{thm:poles-eigenvalues}, \ref{thm:end-to-end}; Section 12.\\
\texttt{TheFractionalRiccatiMuntzPadeTheorem.tm} & Theorem-level preprint; main FRMP theorem; closure under time integration and fractional integration; rough-Heston application; table of coefficient translations. & Thms. \ref{thm:tau}, \ref{thm:cgf}, \ref{thm:global-positive}; Def. \ref{def:cgf-coefficients}; Appendix B.\\
\texttt{residue\_price.tex} & Detailed residue price proof; Lewis contour; intrinsic payoff-pole residues; Padé pole contributions; Lagrange-inversion update discussion; verification table. & Thms. \ref{thm:lewis}, \ref{thm:residue-price}; Section 10.\\
\texttt{single-series-pricing.md} & One-index erfc--Hermite call series; Cayley compactification; coefficient recurrences; absolute convergence; geometric remainder; Black--Scholes certificate; implementation guidance. & Lem. \ref{lem:erfc-kernel}; Thm. \ref{thm:single-series}; Cor. \ref{cor:black-scholes}; Thm. \ref{thm:single-error}; Section 12.\\
\texttt{frmp2.tex} & Concise method sketch; Padé poles $\eta_j$, residues, and rough-Heston evaluation sequence. & Notation section; Thm. \ref{thm:residue-price}; Def. \ref{def:arb4j-flow}.\\
\texttt{TodaLattice.tex} & Background on Flaschka variables, Lax pairs, Jacobi operators, orthogonal polynomials, spectral measures, Riemann--Hilbert view, and Toda hierarchy. & Section 9; Thm. \ref{thm:toda}; Def. \ref{def:jacobi}.\\
\bottomrule
\end{longtable}

\section{Notation translations}

\begin{longtable}{p{0.28\textwidth}p{0.28\textwidth}p{0.34\textwidth}}
\toprule
Source symbol & Canonical symbol & Explanation\\
\midrule
\endfirsthead
\toprule
Source symbol & Canonical symbol & Explanation\\
\midrule
\endhead
$\alpha$ as fractional order in older rough-Heston notes & $\mu$ & $\mu\in(0,1]$ is fixed globally; rough Heston has $\mu=H+\tfrac12$.\\
$t$ and $T$ used interchangeably & $t$ for physical time, $T$ for maturity & $z=t^\mu$ or $z=T^\mu$ according to context.\\
$x=t^\mu$ in some drafts & $z=t^\mu$ & $x$ is reserved for orthogonal-polynomial variables.\\
$a(k,u)$ multiplying $t^{(k+1)\mu}$ & $a_{k+1}(u)$ & Index shift from the TeXmacs preprint.\\
$m_k$ moments in Padé theory & $a_{k+1}$ or $d_k$ & Use $a$ for Riccati moments and $d$ for cgf moments.\\
$P(u),Q(u),R(u)$ in Heston papers & $c_0,c_1,c_2$ in scalar theory & The Heston substitution is $c_0=P(u)$, $c_1=Q(u)$, $c_2=R(u)$.\\
$u$ in pricing contours & $v$ or $w$ & $v$ is the Fourier variable; $w$ is the Lewis contour variable after rotation.\\
$c_j$ as residues in pricing drafts & $r_j$ & Avoids conflict with Riccati coefficients $c_0,c_1,c_2$.\\
$\hat u_j$ pole after rotation & $\eta_j$ or rotated $\hat u_j$ & $\eta_j$ is the canonical Padé pole; hats are used only when rotation is explicit.\\
$\mu_g$ spectral measure & $\nu_g$ & Avoids conflict with fractional order $\mu$.\\
$A_n,B_n$ recurrence coefficients & $\alpha_n,\beta_n$ & Standard Jacobi notation in this exposition.\\
$H_n$ both Hankel determinant and Hermite polynomial & $\tau_n$ for determinant, $H_n^{\mathrm{Herm}}$ when needed & Prevents collision in pricing kernels.\\
$S(n,z)$ generating polynomial & reciprocal orthogonal data & It encodes $z^np_n(1/z)$ and the sigma table.\\
$R_M$ rational approximant & $R_M=P_M/Q_M$ for scalar theory, $\Phi_M$ for cgf & Separates generic Padé notation from pricing notation.\\
$Q_M$ denominator and Heston coefficient $Q(u)$ & $Q_M$ versus $c_1=Q(u)$ & Subscripts distinguish the Padé denominator.\\
$\rho$ as contraction ratio and jump weight & $\rho_*$ for observed ratio, $\omega_g$ for jump weight & Prevents conflict in convergence statements.\\
$\Delta$ as difference and discriminant & $\Delta_M$ for Padé differences, $\Delta_F$ for Fox--Wright discriminant & Context is explicit.\\
\bottomrule
\end{longtable}

\section{Summary of reconciled discrepancies}

\begin{enumerate}[label=(\arabic*)]
\item The indexing of the Müntz coefficients differs across files.
The canonical relation is $a(k)=a_{k+1}$.
\item The Caputo IVP and the Riemann--Liouville Volterra equation are equivalent for the zero-initial-value solution.
The exposition uses Caputo for the differential equation and Riemann--Liouville for coefficient derivations.
\item Local Puiseux convergence is not global series convergence.
Global continuation is supplied by Cole--Hopf meromorphy and Padé convergence theorems.
\item The statement that every Cole--Hopf pole is simple is generic.
The canonical version records the nondegeneracy condition $\Dc w(z_0)\ne0$ at a zero of $w$.
\item The broad ``global convergence'' claims in method drafts require the sharper RH-view hypotheses:
no true pole on the compact set and the Gončar condition for the Padé table.
\item Hankel determinant formulas are canonical algebraically,
but Chebyshev--Wheeler is the canonical computational route.
\item The residue expansion and the one-index erfc--Hermite expansion both evaluate the same Lewis integral under their stated assumptions.
The one-index expansion is preferred for implementation because it uses scalar coefficient recurrences and a direct geometric error bound.
\item The symbol $\mu$ is reserved for fractional order.
Spectral measures use $\nu_g$.
Residues use $r_j$ rather than $c_j$.
\end{enumerate}

\bibliographystyle{plain}
\bibliography{refs}

\end{document}
